\section*{v2.3.0 Mathematical Deepening Core}
\begin{quote}
\textbf{Normalization note.} This block is retained as a historical deepening precursor. Its mathematical thickening remains part of the document's proof memory, but the normalized architecture reads it through the later unfolding phases rather than as a second autonomous mathematical foundation.
\end{quote}

This version continues the return from governance-style consolidation to
explicit mathematical derivation. The focus is no longer merely that closure can
be encoded by a defect functional, but that the induced closure flow can be
separated into invariant, dissipative, and residual channels. In this way the
HC term becomes a mathematically trackable transport law rather than a purely
qualitative stabilizer.

\subsection*{1. Orthogonal projector formulation of closure}
Let $\mathcal H=\mathcal H_A\oplus\mathcal H_B\oplus\mathcal H_C\oplus\mathcal H_R$
be the structural state space used in the previous thickening step, and let
$\Pi_{\mathrm{cl}}$ denote the orthogonal projector onto the admissible closure
manifold
\[
\mathfrak M_{\mathrm{cl}}:=\{\Psi\in\mathcal H\,:\,\Gamma(\Psi)=0\}.
\]
Whenever $\Gamma$ is locally quadratic near $\mathfrak M_{\mathrm{cl}}$, one may
write, to leading order,
\[
\Gamma(\Psi)\asymp \| (I-\Pi_{\mathrm{cl}})\Psi\|^2,
\]
so that the distance to closure is measured by the defect component
\[
\Psi_{\perp}:=(I-\Pi_{\mathrm{cl}})\Psi,
\qquad
\Psi_{\parallel}:=\Pi_{\mathrm{cl}}\Psi.
\]
The compensated flow
\[
\partial_\tau \Psi = F(\Psi)-\kappa\nabla\Gamma(\Psi)
\]
therefore admits the splitting
\[
\partial_\tau \Psi_{\parallel}=\Pi_{\mathrm{cl}}F(\Psi)+\mathcal N_{\parallel}(\Psi),
\qquad
\partial_\tau \Psi_{\perp}=(I-\Pi_{\mathrm{cl}})F(\Psi)-\kappa\nabla\Gamma(\Psi)+\mathcal N_{\perp}(\Psi),
\]
where $\mathcal N_{\parallel},\mathcal N_{\perp}$ collect higher-order curvature
terms of the closure manifold.

\paragraph{Lemma 1 (normal damping near closure).}
Assume that $\Gamma$ is $C^2$ in a neighborhood of $\mathfrak M_{\mathrm{cl}}$ and that
its Hessian is strictly positive on the normal bundle of
$\mathfrak M_{\mathrm{cl}}$. Then there are constants $c_1,c_2>0$ such that for all
states sufficiently near closure,
\[
c_1\|\Psi_{\perp}\|^2\le \Gamma(\Psi)\le c_2\|\Psi_{\perp}\|^2,
\]
and the HC term damps the normal component at leading order.

\begin{proof}
By the Morse--Bott type expansion of $\Gamma$ around the zero set,
\[
\Gamma(\Psi_{\parallel}+\Psi_{\perp})=
\tfrac12\langle H_{\Gamma}(\Psi_{\parallel})\Psi_{\perp},\Psi_{\perp}\rangle
+O(\|\Psi_{\perp}\|^3),
\]
with $H_{\Gamma}$ positive on normal directions. Hence $\Gamma$ is equivalent to
the squared normal distance, and
$\nabla\Gamma(\Psi)=H_{\Gamma}(\Psi_{\parallel})\Psi_{\perp}+O(\|\Psi_{\perp}\|^2)$.
Therefore the compensator contributes a strictly contracting term in the normal
direction.
\end{proof}

\subsection*{2. Differential inequality for the closure ladder}
Set $E(\tau):=\Gamma(\Psi(\tau))$. Then
\[
\dot E = \langle \nabla\Gamma(\Psi),F(\Psi)\rangle - \kappa\|\nabla\Gamma(\Psi)\|^2.
\]
To make the structural content explicit, split the drift as
\[
F = F_{\mathrm{tan}} + F_{\mathrm{nor}} + F_{\mathrm{res}},
\]
where $F_{\mathrm{tan}}$ is tangent to $\mathfrak M_{\mathrm{cl}}$, $F_{\mathrm{nor}}$ is
controlled by the closure defect, and $F_{\mathrm{res}}$ contains genuinely residual
terms. Then
\[
\dot E
= \langle \nabla\Gamma,F_{\mathrm{nor}}\rangle
 + \langle \nabla\Gamma,F_{\mathrm{res}}\rangle
 - \kappa\|\nabla\Gamma\|^2,
\]
because the tangential part drops out to first order.

\paragraph{Proposition 1 (closure ladder with residual threshold).}
Suppose there exist constants $a<\kappa$ and $b\ge 0$ such that
\[
|\langle \nabla\Gamma,F_{\mathrm{nor}}+F_{\mathrm{res}}\rangle|
\le a\|\nabla\Gamma\|^2 + b\Gamma.
\]
Then for states near closure,
\[
\dot E \le -(\kappa-a)\|\nabla\Gamma\|^2 + bE.
\]
In particular, if $b$ is dominated by the coercivity constant relating
$\|\nabla\Gamma\|^2$ and $E$, then the closure ladder is dissipative and the flow
returns exponentially to the admissible structural sector.

\begin{proof}
Insert the estimate directly into the identity for $\dot E$. The final claim
follows from the local coercivity
$\|\nabla\Gamma\|^2\ge cE$ near $\mathfrak M_{\mathrm{cl}}$, which converts the
inequality into $\dot E\le -\mu E$ for some $\mu>0$.
\end{proof}

This proposition sharpens the earlier Lyapunov statement. Closure is not only
monotone under ideal tangency; it remains stable under a controlled class of
residual perturbations, provided the compensator strength exceeds the residual
threshold.

\subsection*{3. Term-wise master equation refinement}
Write the first-layer master equation in the refined form
\[
\partial_\tau \Psi
= T_{\mathrm{tri}}(\Psi)
+ T_{\perp}(\Psi)
+ T_{\mathrm{meas}}(\Psi)
+ T_{\mathrm{res}}(\Psi)
- \kappa\nabla\Gamma(\Psi).
\]
Here the meaning of the terms is:
\begin{itemize}
    \item $T_{\mathrm{tri}}$: internal triolic transport preserving the intrinsic
    algebraic coupling,
    \item $T_{\perp}$: orthogonality-restoring channel induced by the requirement
    that one branch remains perpendicular to the other two,
    \item $T_{\mathrm{meas}}$: frame-shift and measurement transport,
    \item $T_{\mathrm{res}}$: unresolved residual transport not yet assigned to a
    unique physical model.
\end{itemize}
The closure condition can then be tested term by term through the scalar
pairings
\[
\sigma_{\mathrm{tri}}:=\langle \nabla\Gamma,T_{\mathrm{tri}}\rangle,
\quad
\sigma_{\perp}:=\langle \nabla\Gamma,T_{\perp}\rangle,
\quad
\sigma_{\mathrm{meas}}:=\langle \nabla\Gamma,T_{\mathrm{meas}}\rangle,
\quad
\sigma_{\mathrm{res}}:=\langle \nabla\Gamma,T_{\mathrm{res}}\rangle.
\]
Thus
\[
\dot E = \sigma_{\mathrm{tri}}+\sigma_{\perp}+\sigma_{\mathrm{meas}}+\sigma_{\mathrm{res}}
-\kappa\|\nabla\Gamma\|^2.
\]
The mathematical advantage is that every structural claim can now be tied to one
of these four channels. For instance, a perfect intrinsic closure sector is
characterized by $\sigma_{\mathrm{tri}}=0$, while measurement-induced distortion is
captured by the sign and size of $\sigma_{\mathrm{meas}}$.

\paragraph{Corollary 2 (term-wise admissibility test).}
If each non-compensated channel satisfies an estimate of the form
$|\sigma_j|\le a_j\|\nabla\Gamma\|^2+b_jE$ with $\sum_j a_j<\kappa$, then the
master flow is closure-admissible. In this sense the first-layer master
equation acquires a concrete mathematical admissibility criterion instead of a
purely narrative one.

\subsection*{4. Probe reduction with quantitative search radius}
For the probe functional
\[
\Gamma_{\mathrm{probe}}(\Psi_1,\Psi_2;Q)
=\Gamma(\Psi_1)+\Gamma(\Psi_2)+\rho\,\Xi(\Psi_1,\Psi_2;Q),
\]
define the admissible search radius
\[
\mathcal R_\varepsilon(Q):=
\{(\Psi_1,\Psi_2):\Gamma_{\mathrm{probe}}(\Psi_1,\Psi_2;Q)\le \varepsilon\}.
\]
This produces a genuinely quantitative reading of the earlier search-space
language: the probe channel is restricted to a sublevel tube around the
orthogonally compatible sector.

\paragraph{Proposition 3 (quantitative probe reduction).}
Assume the probe drift satisfies the analogue of the closure-ladder estimate.
Then for every sufficiently small $\varepsilon>0$ the set
$\mathcal R_\varepsilon(Q)$ is forward invariant under the compensated probe flow,
provided the initial probe state lies in $\mathcal R_\varepsilon(Q)$. Hence the
probe search does not merely heuristically prefer a smaller space; it is
mathematically confined to a controlled admissible tube.

\begin{proof}
Apply the same differential-inequality argument to
\[
E_{\mathrm{probe}}(\tau):=\Gamma_{\mathrm{probe}}(\Psi_1(\tau),\Psi_2(\tau);Q).
\]
If $\dot E_{\mathrm{probe}}\le -\mu E_{\mathrm{probe}}$ or, more generally, $\dot E_{\mathrm{probe}}\le 0$ on the boundary level set $E_{\mathrm{probe}}=\varepsilon$, trajectories cannot leave the sublevel region.
\end{proof}

\subsection*{5. Consequence for the next development step}
The present deepening changes the status of the HC/closure block in two ways.
First, closure is now expressed through projector splitting and coercive defect
control. Second, the master and probe channels both admit explicit admissibility
tests in terms of scalar defect pairings. The next natural step is therefore no
longer another meta-level consolidation, but the specialization of
$T_{\mathrm{tri}},T_{\perp},T_{\mathrm{meas}},T_{\mathrm{res}}$ and of the probe drift to
more definite structural or physical models. That is the point at which the
current structure-first proof begins to turn into a truly term-resolved
mathematical derivation.


